\section{Finite acquisition for one fixed marginal}
\label{ext:finite-acquisition}
The count argument needs more than an upper bound for the defect. Its
input includes the sign probability, the central gap, the two conditional
variances, and bounds for several signed moments. This appendix explains
how the finite calculation acquires them together. The sample size stays
fixed throughout each application.

Suppose $n\ge5$ is a strict first failure, and take the attained positive
closed-event ratio maximizer $(P,H)$. Then $C=C_n>\CE$, and its same-law
predecessor loss is zero. First acquire support by
Lemma~\ref{ext:direct-gaussian-support-criterion} or one of its separately
proved specializations. This step uses the contact inequality and the
Gaussian CDF, without higher moments. Only afterwards use
\[
 \rho=\E|X|^3,\quad \gamma=\E X^3,\quad
 \beta=\E[X|X|],\quad m=\E X_+=\E X_-,\quad \mu_j=\E X^j.
\]
The signed-moment and Fourier formulas in Appendix~\ref{re:atlas} then
apply to this supported marginal. At small indices their numerical
integrals use the literal value of $n$ and, where needed,
\eqref{ext:finite-low}. No later-index approximation is substituted.

\subsection{What a state asserts}
A state is a finite list of simultaneous inequalities for this one law:
closed intervals for $\rho,\beta,\gamma,m$, signed support radii, a
score ceiling, and any already established raw, defect, or threshold
bounds. Once a central gap is known, it also contains intervals for
$p=\Pp(X>0)$ and its two endpoints. The inequalities are necessary
conditions. Their feasible set may contain points that do not represent
any probability law, which is harmless for an upper bound.

The following rules preserve the actual parameter point. Intersect a
coordinate interval with a newly proved necessary interval; replace an
upper bound by the smaller of two valid upper bounds; or subdivide a
closed interval into two closed children sharing their endpoint.
A negative upper bound for a nonnegative quantity excludes a cell.
A failed denominator, Taylor-domain, or positivity gate does not exclude
the law: that proposed bound is unavailable, and the cell must receive
another valid bound or remain unresolved.

\subsection{Moment contractions with the correct sign radii}
Write $a_\pm=(1\pm\beta)/2$, $b_\pm=(\rho\pm\gamma)/2$.
The inequalities
\[
 m^2\le a_+a_-,\qquad a_\pm^2\le m b_\pm,\qquad
 m\ge\frac{1+\beta^2}{2\rho}
\]
retain the same $m$ on both signs. They also hold when there is a zero
atom. Thus solving them on a closed moment box is a valid contraction,
but choosing a different mean separately for different tests is not an
acquisition of one common state.

Here is the version of the moment matrices that permits different
positive and negative support radii. On either sign write $Y=|X|$,
$0\le Y\le r$, and let $p_s,m,a,b,c_4,c_5$ be its unnormalized
moments of orders zero through five. Positivity of the weights
$r-Y$, $Y(r-Y)$, and $Y^2(r-Y)$ gives
\[
 \begin{pmatrix}rp_s-m&rm-a\\rm-a&ra-b\end{pmatrix},\quad
 \begin{pmatrix}rm-a&ra-b\\ra-b&rb-c_4\end{pmatrix},\quad
 \begin{pmatrix}ra-b&rb-c_4\\rb-c_4&rc_4-c_5\end{pmatrix}
 \succeq0.
\]
For positive denominators their determinant inequalities imply
\begin{align*}
 p_s&\ge m/r+(rm-a)^2/[r(ra-b)],\\
 c_4&\le rb-(ra-b)^2/(rm-a),\\
 c_5&\le rc_4-(rb-c_4)^2/(ra-b).
\end{align*}
A zero diagonal entry forces the corresponding off-diagonal entry to
vanish. One may check that condition or discard the nonnegative
correction, but one cannot divide by zero. In interval form a correction
is bounded below by the square of a positive-part lower numerator
divided by an upper positive denominator. The sign radius $r$ is used
only for its own class. Summing the first inequality over the two signs
and using $p_++p_-\le1$ can improve the common upper mean.

One further bound used at small indices strengthens the signed constant
coefficient. Set $\delta=\gamma-\rho\beta$ and $Y=|X|$. Then
\[
 V:=\E\{Y(Y^2-\rho Y)^2\}
   =\E|X|^5-2\rho\mu_4+\rho^3,
 \qquad
 \delta=\E\{Y\operatorname{sgn}(X)(Y^2-\rho Y)\}.
\]
Here $m>0$ by centering and unit variance.
Cauchy's inequality for the nonnegative measure $Y\,dP$ gives
$\delta^2\le2mV$. Consequently
\begin{equation}\label{ext:weighted-cubic-loss}
 \frac\rho2\mu_4-\frac14\E|X|^5
 \le\frac{\rho^3}{4}-\frac{(\gamma-\rho\beta)^2}{8m}.
\end{equation}
This deduction can be combined with the two other deductions in
Appendix~\ref{re:atlas} by taking their maximum. Their sum would spend
the same nonnegative quantity more than once.

\subsection{The first raw bound comes before defect feedback}
Let $e=C-\CE>0$, $w=H^2/(2n)$ and $\phi_0=\phi(0)$.
The exact conserved identity, with nonnegative contact loss $L$, is
\[
 ne\rho+n\CE\rho(1-e^{-w})+
       n\phi_0e^{-w}D/6+n^{3/2}L=\mathcal R.
\]
Initially use only $D\ge0$, acquired support and moments, and complete
LOW and HIGH estimates. The first full cover must prove a raw upper
bound $\mathcal R\le r<2n/5$. In particular, it must not assume the
defect bound that it proposes to export.

Since $C\rho>\CE>2/5$, rearrangement now gives
\[
 \phi_0e^{-w}(\gamma+3J)/6
       =C\rho+\sqrt nL-\mathcal R/n>0.
\]
This is precisely the sign license needed for removing the Gaussian
weight. Lemma~\ref{re:weight}, used here at $N=n$, therefore proves
\[
 ne\rho+n\phi_0D/6+n^{3/2}L\le\mathcal R,
 \qquad e^{-w}\ge1-5r/(2n),\qquad D\le15000r/(997n).
\]
Every later refinement may use these inequalities, always from a
previous completed pass. The optional additional phase curvature is
\eqref{re:extra}; it is included once in the budget used to update $D$.
The raw remainder and that smaller budget have different meanings.

The signed moment in the actual identity remains $\gamma$, even if an
auxiliary Fourier estimate is evaluated after reflection. In particular,
when $\gamma\ge g_0$ and $\phi\in[\phi_l,\phi_u]$, the valid
curvature lower bound is
\[
 C\rho+\tfrac12\min\{\phi_lg_0,\phi_ug_0\}.
\]
Section~\ref{re:signed-phase-repair} proves the coherent corner rule
for this expression. Complete score and moment corners are taken before
maximizing over $H$ or completing the square. If two complete Fourier
estimates are available, their resulting upper bounds may be minimized;
mixing their separate coefficients is not justified.

\subsection{From the acquired gap to the count input}
With $\sigma=\sqrt{10}$, the bound
$D<\sigma\rho-\gamma$ excludes zero from the support, because otherwise
$J\le\rho$. Both signs have positive probability, so a genuine central
gap $(-A,B)$ follows. The signed endpoint, central-integral and outer-gap
bounds of Appendix~\ref{re:atlas} can now refine the state. In particular,
the complete central rectangle is bounded by its tensor Bernstein
coefficients, using the lower mean for a negative coefficient and the
upper mean for a positive coefficient.

The joint inequality~\eqref{ext:joint-physical-resource} then acquires
the triangle~\eqref{ext:count-triangle-derived}. Its remaining defect
budget is shared by inward displacement and conditional variance.
All refinements of conditional moments price their entire nonnegative
linear cost under this single budget, as in Lemma~\ref{re:pricing}.
This supplies the actual common-law hypotheses for the count calculation
and the optional filters in Appendix~\ref{ext:count-filters}.

The induction is finite. At each complete pass, every actual parameter
point belongs to an initial cell and remains in a child or in a cell
with a proved exclusion. A declared grid with sizes
$(N_\rho,N_\beta,N_\gamma,N_m)$ must account for exactly
$N_\rho N_\beta N_\gamma N_m$ cells; skipping an excluded block must
account for its full multiplicity. An export to the probability panel
$[j/100,(j+1)/100]$ takes maxima of upper bounds and minima of lower
bounds over every active cell touching that panel. In particular,
an endpoint at $j/100$ belongs to both adjacent panels. This is why
the export range is
\[
 \lceil100p_l\rceil-1\ \le j\le\ \lfloor100p_u\rfloor,
\]
intersected with $\{0,\ldots,99\}$. An empty intersection or a strict
necessary-inequality contradiction can close a cell; an incomplete
calculation cannot.

The resulting proof is an induction on finite enclosures, not a
convergence claim about an iterative search. Its last obligations are
concrete: the complete covers must pass, each resource must consume its
proved initializer, and each count contract must be implied by the
resource for the same moment and probability panel. These are the joins
checked by the accompanying reproduction programs.

\subsection{The actual-index acquisitions at $35$ through $39$}
\label{ext:actual3539-acquisition}
These five steps use the actual index in every INITIAL and resource
bound. They therefore need no transfer from a smaller evaluation index.
For $35\le n\le38$, Lemma~\ref{ext:first-failure25} supplies
$C_n<41/100+4177/(13824n)$. This is below all four native INITIAL
ceilings in Section~\ref{ext:finite1539} and also below $423493/10^6$, the smaller
ceiling used by the local calculations at $35$.
The sharp predecessor gives contact allowance zero; the weaker allowances
$19/100$ and $63/340$ retained in the index-$35$ support and resource
programs are consequently valid. At $39$, use the separately proved
positive-allowance step to obtain the smaller ceiling $41449/10^5$.

The global INITIAL covers have respectively $10,16,12,15,12$ finite
refinement passes, followed by a full $100$-panel export. Each pass
retains one law's earlier bounds, proves support before using its
higher moments, and rounds upper bounds upwards and lower bounds
downwards. The global support bounds are $39/10$ at $35$ and $19/5$
at $36$--$39$. At $35$, the separate support calculation actually gives
$96/25<39/10$. Local refinements use separately checked smaller radii.

At $35$, four short INITIAL refinements cover the intervals
$[.29,.34]$, $[.34,.39]$, $[.39,.45]$ and $[.45,.57]$;
they use respectively $4,3,3,2$ passes. Bounds shared by a whole
interval are the weakest bounds over its still possible panels:
the maximum of their upper bounds and the minimum of their lower
bounds. These are intersected with the already acquired bounds for
that interval. Thus no moment from one law is assigned to another.
The resulting count interfaces cover $P_{31},\ldots,P_{48}$.
Three further alternating INITIAL/resource calculations cover
$P_{16},\ldots,P_{28}$: they exclude $P_{16},\ldots,P_{22}$
and produce count interfaces on $P_{23},\ldots,P_{28}$.
The first short interval also supplies the two interfaces on
$P_{29},P_{30}$. Finally the global source, supplemented by the last
short interval on $P_{49},\ldots,P_{56}$, excludes every panel in
$P_0,\ldots,P_{15}$ and $P_{49},\ldots,P_{99}$.
The union consists of $74$ excluded panels and $26$ count panels,
with every probability endpoint included.

At $36$, a two-pass localized INITIAL calculation on $[.30,.42]$
replaces the global source only on $P_{30},\ldots,P_{41}$; its proved
one-sided support radii are $3$ and $37/10$. All other panels retain
their global source. At $37$--$39$, the global exports are used directly.
The full resource traversals leave respectively
$P_{24},\ldots,P_{48}$, $P_{24},\ldots,P_{47}$,
$P_{25},\ldots,P_{47}$ and $P_{26},\ldots,P_{47}$.
Every other panel reaches an explicit moment, support or budget
contradiction during the finite traversal.


At $35$ the phase budget also retains the actual excess. Since
$C>\CE$ and $\CE<20487/50000$, its lower charge is
$35\max\{0,C-20487/50000\}\rho$. For each fixed phase, the resulting
upper budget is affine in $C$ on either side of this breakpoint,
and affine in $\rho$ for fixed $C$. Its maximum on the score--moment
rectangle therefore occurs among the six corners obtained from the
two score endpoints, the breakpoint, and the two moment endpoints.
Subtracting this charge before the phase comparison is justified by
the same conserved identity used in Appendix~\ref{ext:finite-floor-transfer}.

Each surviving row supplies upper bounds for the gap, noise,
amplitudes and the two signed phases, together with one simultaneous
affine defect lower bound. The count contract weakens each of these
bounds in the valid direction. For the phase envelopes, every acquired
corner polynomial is bounded above on the whole signed interval by
at least one count-certificate corner polynomial. The comparison checks
the endpoints and, when it lies inside the interval and the difference
is convex, its vertex. Taking maxima therefore preserves the
inequality. These comparisons connect the freshly acquired same-law
rows to the $26+25+24+23+22=120$ actual-index count forests.
